\clearpage
\begingroup
\makeatletter
\begingroup\lccode`+=32 \lowercase
 {\endgroup\def\Url@ObeySp{\Url@Edit\Url@String{ }{+}}}
\def\Url@space{\penalty\Url@sppen\ }
\makeatother
\section{The genus-two partition-function experiment}
\label{app:partition-test}

This appendix documents the saved September 15--16, 2026 comparison
between an NS--R--R theta decomposition and an all-NS decomposition of
the same marked surface. It supplements the nonchiral conventions in
Appendix~\ref{app:nonchiral-conventions}; it is not an additional
chiral-block identity. With the historical sewing matrix
\begin{equation}
 M_{fg,\mathrm{old}}^{(\eta,\eta')}=
 \frac{B_{L,\eta}B_{R,\eta'}}8
 \begin{pmatrix}1&\eta_1\eta_3\\\eta_1\eta_3&1\end{pmatrix}_{fg},
 \qquad \eta\eta'=-\eta_2\eta_3,
 \label{eq:partition-historical-matrix}
\end{equation}
the normalized cross-channel ratio converges near $1/4$.
A subsequent experiment multiplies that entire
matrix by four, as an explicit assumption. Its remaining discrepancies
are below $0.1\%$ at the central tested surface. The calculation does
not derive this multiplier or establish the proposed interacting spin
transport. These historical normalization claims are superseded by
Sec.~\ref{app:partition-normalization-resolution}: the actual Liouville
identity residue fixes the missing factor two at each NS--R--R vertex.
The current matrix is \eqref{eq:app-nonchiral-matrix}. All tables before
that final subsection retain their original normalization.

The main record is
\path{Machine Notes/Genus 2/NSRR_BILINEAR_PAIRING_2026-09-15.md}.
The completed numerical reports are
\path{Data Set/nsrr_bilinear_cross_channel_20260915/README.md},
\path{Data Set/nsrr_bilinear_quadrature_20260915/FINAL_REPORT.md}, and
\path{Data Set/nsrr_provisional_factor4_20260916/README.md}.
The latter two supersede the central low-order estimates in the first.
The tables below transcribe those archived computations; adding this
appendix does not rerun their momentum integrals with the subsequently
revised release code.

\subsection{Observable, spectrum, and primary propagation}

The continuum input is delta-normalized $\mathcal N=1$ super-Liouville
theory, with
\begin{equation}
 b=1.4,\qquad Q=b+b^{-1},\qquad
 c=\frac32+3Q^2,\qquad
 h_{\NS}(p)=\frac{Q^2}{8}+\frac{p^2}{2},\qquad
 h_{\R}(p)=h_{\NS}(p)+\frac1{16}.
 \label{eq:partition-spectrum}
\end{equation}
Thus $P=ip$ in the chiral momentum convention of these notes, and
each internal continuum label is integrated over $p>0$ with measure
$dp/\pi$. Both channels omit the same momentum-independent
cosmological prefactor; the saved settings are
\texttt{mu=1+0j} and \texttt{include\_cosmological\_prefactor=false}.

The comparison removes the change of conformal frame using a physical
free scalar plus one real Majorana fermion, of central charge $3/2$:
\begin{equation}
 \mathcal R=
 \frac{Z_{\NS\R\R}}{Z_{\NS\NS\NS}}
 \left(\frac{Z_{\mathrm{free,target}}}
             {Z_{\mathrm{free,source}}}\right)^{\!\kappa},
 \qquad \kappa=\frac{c}{3/2}=1+2Q^2
              =9.940408163265307.
 \label{eq:partition-ratio}
\end{equation}
Agreement in the tested prescription requires $\mathcal R=1$.
The free theory here is the reference used to remove the conformal
anomaly; it is distinct from the auxiliary fermion in the
double-Virasoro convolution. Raw partition functions in the two
plumbing frames need not be equal.

The data files order the edges geometrically as $(0,1,\infty)$,
namely $(\R_0,\R_1,\NS)$ for the source. Reverse this order when
calling a block in the $(\NS_\infty,\R_1,\R_0)$ slots of these notes.
For example, source primary propagation is
\begin{equation}
 \exp\!\left[h_{\R}(p_0)\operatorname{Log} q_0+
 h_{\R}(p_1)\operatorname{Log} q_1+h_{\NS}(p_\infty)\operatorname{Log} q_\infty\right].
 \label{eq:partition-primary}
\end{equation}
The target uses its own three plumbing parameters and three NS
weights. These factors multiply the descendant blocks from outside;
neither the blocks nor the nonchiral matrix contain an additional
$q^h$ or cylinder factor $q^{-c/24}$. The archived integration uses
the principal logarithms recorded at each node. More general
continued logarithms and independent antichiral weights were checked
separately. Complex conjugation is imposed only on the physical real
slice used for the integrals.

\paragraph{Three-point input.}
The supplied constants are the real-continuum coefficients
$(C,\widetilde C)$ for NS--NS--NS and $(E,O)$ for R--R--NS,
in the common two-point normalization $\pi\delta(p-p')$.
The implementation is
\path{Code/c_Recursion/generic_super_liouville_structure_constants.py},
with calls \texttt{ns\_constants(p0,p1,pinf)} and
\texttt{rr\_ns\_constants(p1,p0,pinf)}. Its definitions use
\begin{align}
 \Upsilon_{\NS}(x)&=\Upsilon_b(x/2)\Upsilon_b((x+Q)/2),\nonumber\\
 \Upsilon_{\R}(x)&=\Upsilon_b((x+b)/2)\Upsilon_b((x+b^{-1})/2).
 \label{eq:partition-upsilon}
\end{align}
Here $\Upsilon_b(Q/2)=1$ and
$\Upsilon'_{\NS}(0)=\Upsilon_b(b)/2$.
For clarity, the leg factors in this implementation are
\begin{align}
 \mathcal N_{\NS}(p)&=
 \sqrt{\Upsilon_{\NS}(2ip)\Upsilon_{\NS}(-2ip)}\,
 |\gamma(1-ip/b)\gamma(-ipb)|^{1/2},\nonumber\\
 \mathcal N_{\R}(p)&=
 \sqrt{\Upsilon_{\R}(2ip)\Upsilon_{\R}(-2ip)}\,
 |\gamma(\tfrac12+ipb)\gamma(\tfrac12+ip/b)|^{-1/2},
 \label{eq:partition-legs}
\end{align}
where $\gamma(x)=\Gamma(x)/\Gamma(1-x)$ and each square root is
positive for positive real momentum. To display the constants
without four repeated long arguments, set
$x_0=Q/2+i(p_1+p_2+p_3)$ and
$x_j=Q/2+i(p_1+p_2+p_3-2p_j)$, $j=1,2,3$, in this paragraph only.
Then the implemented input, with the cosmological factor omitted, is
\begin{align}
 C&=\frac{\Upsilon'_{\NS}(0)\prod_{j=1}^3\mathcal N_{\NS}(p_j)}
                {\prod_{j=0}^3\Upsilon_{\NS}(x_j)},&
 \widetilde C&=\frac{2\Upsilon'_{\NS}(0)\prod_{j=1}^3\mathcal N_{\NS}(p_j)}
                {\prod_{j=0}^3\Upsilon_{\R}(x_j)},\nonumber\\
 E&=\frac{b^{-2}\Upsilon'_{\NS}(0)
          \mathcal N_{\R}(p_1)\mathcal N_{\R}(p_2)\mathcal N_{\NS}(p_3)}
 {\Upsilon_{\R}(x_0)\Upsilon_{\R}(x_3)
  \Upsilon_{\NS}(x_1)\Upsilon_{\NS}(x_2)},&
 O&=\frac{b^{-2}\Upsilon'_{\NS}(0)
          \mathcal N_{\R}(p_1)\mathcal N_{\R}(p_2)\mathcal N_{\NS}(p_3)}
 {\Upsilon_{\NS}(x_0)\Upsilon_{\NS}(x_3)
  \Upsilon_{\R}(x_1)\Upsilon_{\R}(x_2)}.
 \label{eq:partition-constants}
\end{align}
The Ramond leg metrics and the relative $b^{-2}$ must be retained at
$b=1.4$; their self-dual values do not fix the generic-$b$ convention.
The odd NS vertex has the phase $i\widetilde C$ in the Human Note
convention. That phase is handled by sewing, not inserted a second
time into the real coefficient returned by the constants routine.

\subsection{Local nonchiral sewing and its independent checks}

Equation \eqref{eq:app-nonchiral-projection} defines the projected
physical chiral blocks, and \eqref{eq:partition-historical-matrix} supplies
their local matrix. The historical note denotes these projected
blocks by $\widehat F$. That hat does \emph{not} mean the enlarged
SCA--fermion block. We retain the explicit subscript ``projected''
here to avoid confusing the two constructions.

Write $s=\eta_1\eta_3$ and $r=\eta_2\eta_3$ for this experiment,
where the $\eta_i$ are physical tube signs in slot order
$(\NS_\infty,\R_1,\R_0)$. Simultaneously reversing all three
leaves the full-state tube character unchanged. The local matrix
is diagonal in the vertex-sign pair $(\eta,\eta')$ and supports
$\eta\eta'=-r$. Consequently $r=-1$ selects $E_LE_R$ and
$O_LO_R$, whereas $r=+1$ selects $E_LO_R$ and $O_LE_R$.
These physical tube signs are separate from the formal-parity
evaluations in the projected-block definition.

\paragraph{The physical Ramond dual.}
The finite Clifford calculation explains both the inverse Gram and
the factor $1/8$. Put $u=e^{-i\pi/4}$ and define the odd map
\begin{equation}
 J=\begin{pmatrix}0&u\\\bar u&0\end{pmatrix},\qquad
 J(Ww)=(-1)^{\#G_W}WJw.
 \label{eq:partition-J}
\end{equation}
Here $W$ is a PBW word and $\#G_W$ counts its $G$ modes.
It squares to one, commutes with $L$, anticommutes with $G$, and
reverses the bilinear pairing: $B(Jx,Jy)=-B(x,y)$.
At each level, an even basis $v_i$ and its partners $Jv_i$ separate
the descendant multiplicities from a two-dimensional Clifford
factor. Denote the Gram matrix on the even basis by $B_0$.
The antichiral construction uses conjugate zero-mode phases and
the ground metric $\operatorname{diag}(1,-i)$.

In chiral/antichiral Clifford order $(00,01,10,11)$, the graded
metric and the physical ket and dual embeddings are
\begin{equation}
 D_J=\operatorname{diag}(1,-1,-1,-1),\qquad
 E_R^J=\frac1{\sqrt2}
 \begin{pmatrix}1&0\\0&\bar u\\0&u\\-i&0\end{pmatrix},\qquad
 E_L^J=\frac1{\sqrt2}
 \begin{pmatrix}1&0\\0&-u\\0&-\bar u\\-i&0\end{pmatrix}.
 \label{eq:partition-embeddings}
\end{equation}
Direct multiplication gives $(E_L^J)^TD_JE_R^J=I_2$.
The physical multiplicity Gram is therefore
$B_{0,h}\otimes B_{0,\tilde h}\otimes I_2$.
If the columns of $U$ express these states in the canonical physical
PBW basis, the inverse Gram in that basis is
$U(B_{0,h}^{-1}\otimes B_{0,\tilde h}^{-1}\otimes I_2)U^T$.
The transpose is bilinear, without complex conjugation.
Both trinion tensors take physical \emph{ket} states; $E_L^J$
enters the dual contraction, not a second ket at a vertex.

An independent two-point Ward calculation uses the identity NS
insertion and the coordinate map $z\mapsto z/(1-z)$. For left and
right states $l,r$ it gives
\begin{equation}
 B(l,r)=(-1)^{N_l+\widetilde N_l}
 T_{\mathbf1}\!\left(e^{L_1+\widetilde L_1}l,
                     e^{-L_1-\widetilde L_1}r\right),
 \label{eq:partition-identity-pairing}
\end{equation}
where $N_l,\widetilde N_l$ are the two descendant levels and
$T_{\mathbf1}$ is the full three-point form with that identity
insertion. The historical note reports 228 exact matrix comparisons at
Ramond level pairs $(0,0),(1,0),(0,1),(1,1),(2,0),(0,2)$;
all agree exactly in that record.

\paragraph{Reduction of the two Ramond families.}
For NS chiral and antichiral parities $a,b$, and Ramond
chiral/antichiral parities $(\alpha,\beta)$ and
$(\gamma,\delta)$, separating the two algebras at a vertex gives
$(-1)^{b(a+\alpha+\gamma)+\beta\gamma}$. The $ab$ term comes
from the outgoing NS bra. After removing the multiplicity
amplitudes, the physical-family matrices are
\begin{center}
\begin{tabular}{ccccc}
\toprule
$(a,b)$ & $(0,0)$ & $(0,1)$ & $(1,0)$ & $(1,1)$\\
\midrule
matrix & $\begin{pmatrix}1&0\\0&\eta\end{pmatrix}$
& $\begin{pmatrix}0&\bar u\\\eta u&0\end{pmatrix}$
& $\begin{pmatrix}0&u\\\eta\bar u&0\end{pmatrix}$
& $\begin{pmatrix}-i&0\\0&i\eta\end{pmatrix}$\\
\bottomrule
\end{tabular}
\end{center}
Writing the physical Ramond family parities as
$\epsilon_2,\epsilon_3$, the theta sign is
$K(a+b,\epsilon_2,\epsilon_3)
=(a+b)(\epsilon_2+\epsilon_3)+\epsilon_2\epsilon_3$.
The graded NS inverse Gram contributes $(-1)^{ab}$.
Summing the families and their tube characters gives zero unless
$\eta\eta'=-r$. In the supported sector its values, in the order
$(a,b)=(0,0),(0,1),(1,0),(1,1)$, are
\begin{equation}
 (2,-2is,2is,2).
 \label{eq:partition-family-sum}
\end{equation}
Coefficient by coefficient the projected chiral blocks satisfy
$\mathbf F_1|_{\mathrm{projected}}
=-i(-1)^a\mathbf F_0|_{\mathrm{projected}}$;
the antichiral relation has $+i(-1)^b$.
Each projected even block carries $\sqrt2$. Expressing
\eqref{eq:partition-family-sum} in that basis gives
$\tfrac12\bigl(\begin{smallmatrix}1&s\\s&1\end{smallmatrix}\bigr)$.
The historical adapter took the chiral three-form coefficients to be
$E/2,O/2$, supplying the remaining $1/4$ in
\eqref{eq:partition-historical-matrix}. The Clifford reduction is exact;
this last identification of its input coefficients was the error.

For $r=-1$, before primary propagation, the ground coefficient is
$(E_LE_R+O_LO_R)/2$. With only a holomorphic NS half-level excited,
the coefficient is
\begin{equation}
 is\,\frac{E_LE_R(p_0-p_1)^2+O_LO_R(p_0+p_1)^2}
                {8h_{\NS}},\qquad \beta_j=ip_j/\sqrt2.
 \label{eq:partition-half-level}
\end{equation}
The antichiral half-level has the opposite phase. Their cancellation
on a real slice does not make the separate coefficients zero.

\paragraph{Identity and descendant tests.}
The module test assigned $E=2,O=0$, giving ordinary Ramond trace two and
supertrace zero. After converting the two-point local coordinates,
the ordinary character is
\begin{equation}
 2\prod_{n\geq1}\frac{1+q^n}{1-q^n}
  \prod_{n\geq1}\frac{1+\widetilde q^n}{1-\widetilde q^n}.
 \label{eq:partition-character}
\end{equation}
The checked multiplicities per family are $1,2,4,8$ through level
three; the supertrace vanishes at each tested pair of levels.
This check already counts both Ramond grounds. It does not supply
an additional factor two per Ramond edge, nor identify a product
of finite-neck plumbing parameters with the torus nome without
the coordinate conversion.

The independent full-state Ward/BPW oracle also checks 1,024
descendant coefficients, including NS level $5/2$, Ramond level
two, mixed descendants, all four tube-sign choices, all four
products of pants coefficients, and independent complex antichiral
parameters. Its maximum scaled discrepancy is $1.477\times10^{-15}$.
Twelve finite-series tests with independent complex moduli,
unequal primary weights, and continued logarithms have norm error
at most $2.221\times10^{-16}$ and phase error at most
$1.995\times10^{-16}$ radians. All 243 nodes in the complete
saved L3/N5 and L5/N3 complex-block data were recombined for every
tube-sign choice; their explicit primary factors agree exactly
at saved precision. These are local checks, not global crossing
tests.

\subsection{Marked surfaces, spin transport, and the free reference}

The five surfaces start from the original marking
\begin{equation}
 \Omega(t)=\begin{pmatrix}i&t+i/2\\t+i/2&i\end{pmatrix},
 \qquad t=0.52,\ 0.56,\ 0.60,\ 0.64,\ 0.68.
 \label{eq:partition-surfaces}
\end{equation}
The source and target use different re-marked theta charts. The
saved composed symplectic matrix from source to target, in the
usual $(A,B;C,D)$ block order, is
\begin{equation}
 \begin{pmatrix}
 -1&0&0&0\\1&-1&0&-2\\0&0&-1&-1\\-1&1&0&1
 \end{pmatrix}.
 \label{eq:partition-symplectic}
\end{equation}
The charge periods obtained independently from charged-boson
sewing obey $\Omega_{\mathrm{charge}}=\Omega_{\mathrm{marked}}+B$,
with
\begin{equation}
 B_{\mathrm{source}}=\begin{pmatrix}0&0\\0&1\end{pmatrix},
 \qquad B_{\mathrm{target}}=\begin{pmatrix}-1&-1\\-1&0\end{pmatrix}.
 \label{eq:partition-period-branches}
\end{equation}
Here $B$ is an integer period shift, unrelated to any three-point
coefficient. For binary characteristics $[\alpha|\beta]$, the
charge characteristic is
$\beta_{\mathrm{charge}}=\beta-B\alpha+\operatorname{diag}B$
modulo two. Reducing its components to zero or one also fixes the
theta translation phase
$\exp\{i\pi[\alpha^TB\alpha/4+
\alpha^T(\beta_{\mathrm{charge}}-\beta)/2]\}$.

The two compared components have the following dictionary.
The raw lifts in the last column are in geometry order, with the
redundant infinity lift fixed to $+1$.
\begin{center}
\begin{tabular}{cccc}
\toprule
Source marked & Target marked & Target charge & Target raw lifts\\
\midrule
$[11|00]$ & $[00|00]$ & $[00|10]$ & $(-,+,+)$\\
$[11|11]$ & $[00|10]$ & $[00|00]$ & $(+,+,+)$\\
\bottomrule
\end{tabular}
\end{center}
The remaining source characteristics $[11|01]$ and $[11|10]$
map to $[01|01]$ and $[01|11]$, respectively. Their target
decomposition contains a Ramond edge, so this all-NS comparison
does not test them. In particular, the integrated source uses
$r=-1$, hence equal vertex signs $\eta=\eta'$. It is not a
partition-function test of the inserted recovery for
$\eta\eta'=-1$.

At $t=0.60$ the actual plumbing inputs, rounded here to 12 digits,
are
\begin{center}
\begin{tabular}{ccc}
\toprule
Edge & Source $q$ & Target $q$\\
\midrule
$0$ & $-0.039389297943-0.025083391996i$
    & $ 0.034728649739-0.025129048103i$\\
$1$ & $-0.040592698060+0.029788087392i$
    & $-0.024917675595-0.001472871220i$\\
$\infty$ & $-0.035159173395+0.025344924433i$
    & $-0.080559050182-0.000078089260i$\\
\bottomrule
\end{tabular}
\end{center}
Use \texttt{config.json} in the cross-channel data directory for
the full stored precision and all five surfaces; do not regenerate
rounded plumbing inputs from this table. The geometry provenance is
also retained in
\path{Data Set/fixed_spin_free_NSrr_20260830/summary.json}.

\paragraph{Computing the free frame.}
In the charge frame, the scalar chiral block has the form
$P(q)\exp(i\pi a^T\Omega_{\mathrm{charge}}a)$, where $a$ is
the two-component loop charge and $P(q)$ is its zero-charge
oscillator factor. This paragraph's $P(q)$ is not the internal
SCA momentum or the primary factor \eqref{eq:partition-primary}.
Charge conservation fixes the infinity charge to $-a_0-a_1$.
Two real fermions give the integer charge lattice, shifted by
$\alpha/2$ in characteristic $[\alpha|\beta]$. Thus one
Majorana chiral factor is the continued square root of
$P(q)\theta[\alpha|\beta](\Omega_{\mathrm{charge}})$, and
\begin{equation}
 Z_{\mathrm{free}}=
 \frac{|P(q)|^3\,|\theta[\alpha|\beta](\Omega_{\mathrm{charge}})|}
      {\sqrt{\det(2\operatorname{Im}\Omega_{\mathrm{charge}})}}.
 \label{eq:partition-free}
\end{equation}
The scalar has $h(a)=a^2/2$, measure $da_0\,da_1$, and connected
target-space zero-mode volume one. These choices fix the two in
the determinant. The nonchiral Majorana factor is the modulus
of the Dirac chiral factor, not its modulus squared.

The implementation
\path{Code/genus_2/fixed_spin_free_plumbing.py} uses oscillator
cutoff 32 for the recorded spin transport, charge-lattice cutoff
five, and 32 radial continuation steps from $q_e\mapsto10^{-6}q_e$
to $q_e$. It fixes the root at positive normalized degeneration
coefficient (NS ground one; two-Ramond ground
$\sqrt2\exp[(\operatorname{Log} q_0+\operatorname{Log} q_1)/16]$).
At the central surface the same-frame free factors are
\begin{equation}
 Z_{\mathrm{free,source}}=0.5754095923717206,\qquad
 Z_{\mathrm{free,target}}=0.5638924570404817.
 \label{eq:partition-free-values}
\end{equation}
Their ratio is $1.020424347209\ldots$ and its $\kappa$th power
is $1.222601313543\ldots$.
The charge-period residuals at this point are
$2.53\times10^{-13}$ and $6.33\times10^{-11}$;
the independent free modular/frame residual is
$1.72\times10^{-11}$. Across the five surfaces the period
residual is at most $1.31\times10^{-9}$.

The two independently continued source/target Majorana ratios
have the same phase. At the central surface their quotient is
$1.000000000004227-3.7558\times10^{-12}i$; the maximum discrepancy
from one across all five surfaces is $1.020\times10^{-10}$.
The interacting comparison therefore uses relative phase $+1$,
fixed by this free control rather than fitted to an integral.
Extending this free-spin dictionary to the interacting off-diagonal
pairing is a proposal tested here, not a derivation of its fusion
kernel.

\paragraph{Why the full complex lift vector is retained.}
At $c=3/2$, with charge conservation, independent mode sewing
identifies the two free spin blocks with the projected physical
blocks times their primaries, with relative factor $i$ on
$\mathbf F_1|_{\mathrm{projected}}$.
Thus the physical tube-sign combination corresponds to
$(D_{00}-isD_{11})/\sqrt2$, where $D_{00},D_{11}$ denote the two
raw spin components in the source charge frame. It is generally
not one theta square root. At total descendant level three, two
complex plumbing scales and both allowed charge choices give
maximum norm-ratio error $3.646\times10^{-8}$ and phase error
$5.320\times10^{-9}$ radians. Direct Majorana mode sewing
agrees with the independently continued bosonized reference to
$1.888\times10^{-14}$.

For the all-NS target, the four literal lift blocks must first be
converted to the reflected frame. For even NS primaries the
conversion in three-form sector $a=0,1$ is
\begin{equation}
 \mathbf F^{\mathrm{reflected}}_a(\eta_0,\eta_1,1)
 =\frac14\sum_{\lambda_0,\lambda_1=\pm1}
 \sum_{\substack{p_0,p_1,p_\infty\in\{0,1\}\\
                  p_0+p_1+p_\infty=a\ (\mathrm{mod}\ 2)}}
 (-1)^{K(p_0,p_1,p_\infty)}
 (\eta_0\lambda_0)^{p_0}(\eta_1\lambda_1)^{p_1}
 \mathbf F_a(\lambda_0,\lambda_1,1).
 \label{eq:partition-ns-transform}
\end{equation}
The $p_e$ in this finite sum are parity bits, not continuum
momenta. The formula is the finite parity Fourier transform
implemented in \path{Code/genus_2/all_ns_reflected_sewing.py}.
One selects the two target raw lifts in the table only after this
conversion, and then forms the same normalized two-component
combination. Both NS three-form sectors and their supplied
constants contribute. Taking absolute values of individual
literal lift blocks before this step would lose the interference.

The computation retains a $2\times2$ integrated spin matrix
(called \texttt{density} in the files). Explicitly, on the real
slice used for integration its source value is
\begin{align}
 (Z_{ij})_{\mathrm{source}}
 ={}&\int\prod_e\frac{dp_e}{\pi}\,
 \left|\exp\!\sum_e h_e\operatorname{Log}q_e\right|^2
 \sum_{\eta=\pm1}\frac{B_{L,\eta}B_{R,\eta}}4\nonumber\\[-2pt]
 &\quad\times
 \left.\begin{pmatrix}\mathbf F_0^{(\eta,\eta)}\\
                      i\mathbf F_1^{(\eta,\eta)}\end{pmatrix}
       \right|_{\mathrm{projected}}
 \left(\left.\begin{pmatrix}\mathbf F_0^{(\eta,\eta)}\\
                      i\mathbf F_1^{(\eta,\eta)}\end{pmatrix}
       \right|_{\mathrm{projected}}\right)^{\!\dagger}.
 \label{eq:partition-source-density}
\end{align}
Here $i,j\in\{00,11\}$ label the two spin components, and
$B_{L,+}=B_{R,+}=E$, $B_{L,-}=B_{R,-}=O$ in the numerical run.
The dagger in this outer product is ordinary complex conjugate
transpose on this real slice, not the bilinear BPZ operation.
The target matrix uses its own momenta, primary factor and
plumbing parameters:
\begin{align}
 (Z_{ij})_{\mathrm{target}}
 ={}&\int\prod_e\frac{dp_e}{\pi}\,
 \left|\exp\!\sum_e h_{\NS}(p_e)\operatorname{Log}q_e\right|^2
 \sum_{a=0}^1 C_a^2\nonumber\\[-2pt]
 &\quad\times
 \begin{pmatrix}\mathbf F_a^{\mathrm{reflected}}(-,+,+)\\
                 \mathbf F_a^{\mathrm{reflected}}(+,+,+)\end{pmatrix}
 \begin{pmatrix}\mathbf F_a^{\mathrm{reflected}}(-,+,+)\\
                 \mathbf F_a^{\mathrm{reflected}}(+,+,+)\end{pmatrix}^{\!\dagger},
 \qquad C_0=C,\quad C_1=\widetilde C.
 \label{eq:partition-target-density}
\end{align}
The local NS reflection conversion already combines the odd
sewing sign with the two odd-vertex phases; no extra minus
multiplies the $\widetilde C^2$ term here.
The five reported scalar
observables are its quadratic forms on
$(1,-i)/\sqrt2$, $(1,i)/\sqrt2$, $(1,0)$, $(0,1)$, and its
trace, denoted $s=+1$, $s=-1$, resolved 00, resolved 11, and
spin sum. These are row vectors, contracted as
$v(Z_{ij})\overline v^{\,T}$ in the code. The free-frame
ratio is common to both transported components, so
\eqref{eq:partition-ratio} applies to all five observables.
In particular,
\begin{equation}
 Z_{s=+1}+Z_{s=-1}=Z_{00}+Z_{11}.
 \label{eq:partition-spin-sum}
\end{equation}
The trace has no cross-spin interference. Its discrepancy cannot
be repaired by changing only the relative Majorana phase.

\subsection{Numerical design and the unscaled result}

The source cutoff $L$ is total descendant level. The target
parameter $R$ is the twice-level cutoff on NS $c$-recursion
corrections; its large-$c$ global contribution is resummed.
These are different truncations, so R12 is not a polynomial
block truncated in every edge at level 12. The source uses the
native double-Virasoro engine; the target uses
\texttt{NSGenus2CRecursion.collision\_aware\_block\_mp} with
\texttt{global\_method=resummed}. The target vacuum settings are
Schottky word length seven and mode cutoff 50. The relative
global-sum tolerance is $2\times10^{-9}$; its maximum total
occupation was raised from 36 to 60 in the quadrature refinement.
All accepted target nodes report converged global sums.

The source native engine and structure constants use 40 decimal
digits; target recursion uses 50. Stored projected amplitudes,
spin matrices, and quadrature accumulation use complex binary64.
Working precision is therefore not a claim of 40 accurate digits
in the final partition function. The source coefficient banks
retain the complex channels before nonchiral contraction and
are reused at different plumbing parameters.

For the unscaled refinement, each edge uses generalized
Gauss--Laguerre nodes $x_j$ and weights $w_j$ for weight
$x^{-1/2}e^{-x}$ on $[0,\infty)$. With the fixed envelope
$|q_{e,\mathrm{env}}|$, the actual momentum rule is
\begin{equation}
 p_{e,j}=\sqrt{\frac{x_j}{-\log|q_{e,\mathrm{env}}|}},\qquad
 W_{e,j}=\frac{w_j e^{x_j}}
                  {2\pi\sqrt{-\log|q_{e,\mathrm{env}}|}},\qquad
 \int_0^\infty\frac{dp_e}{\pi}\,f(p_e)
       \simeq\sum_{j=1}^N W_{e,j}f(p_{e,j}).
 \label{eq:partition-quadrature}
\end{equation}
The three-dimensional weights are products of these one-edge
weights; the Gaussian primary propagation remains in the
integrand. In geometry order the source envelopes are
$(0.04779213238222936,0.0520035772489429,0.04349541543013214)$,
and the target envelopes are
$(0.04677899744906184,0.03806246230229872,0.1332848304670786)$.
Only complete $N^3$ tensor grids enter any reported integral.

\paragraph{Initial five-surface comparison.}
With source L5, target R12, and N3 in both channels, the local
matrix gives the following ratios. No factor four is applied.
\begin{center}
\begin{tabular}{rrrr}
\toprule
$t$ & $\mathcal R(s=+1)$ & $\mathcal R(s=-1)$ & $\mathcal R(\text{spin sum})$\\
\midrule
0.52 & 0.251461747 & 0.251166853 & 0.251308055\\
0.56 & 0.252437897 & 0.252246849 & 0.252338372\\
0.60 & 0.252925745 & 0.252836380 & 0.252879225\\
0.64 & 0.252916112 & 0.252965657 & 0.252941875\\
0.68 & 0.252456067 & 0.252714296 & 0.252590118\\
\bottomrule
\end{tabular}
\end{center}
At $t=0.60$, changing L3 to L5 at N3 changes the two sign ratios
and the spin-sum ratio by $2.650\times10^{-4}$,
$1.287\times10^{-4}$, and $1.941\times10^{-4}$ relatively.
Changing target R8 to R12 at N3 changes those ratios by
$-3.915\times10^{-7}$, $-1.910\times10^{-8}$, and
$-1.976\times10^{-7}$. Changing both grids from N3 to N4 at
L3/R12 gives $-2.447\times10^{-3}$, $-2.776\times10^{-3}$,
and $-2.618\times10^{-3}$; the resulting spin-sum ratio is
$0.252168291$. The grid effects warrant refinement but are much
smaller than the discrepancy from one.

This stage computes 91 fresh target momentum nodes: 27 shared
nodes evaluated at all five surfaces and 64 at the central
surface. Each retains both form sectors, all four complex
lifts, and both target recursion cutoffs. Independent
reconstruction of 334 source/target measures gives maximum
relative error $6.661\times10^{-16}$.
There are 120 source-block spot comparisons through L5 against
the native engine. Their largest scaled difference,
$8.117\times10^{-8}$, occurs in an extreme tail; the largest
single checked-node correction divided by the integral is
$3.727\times10^{-12}$. This is a three-node spot check, not
a recomputation of the complete source grid. The saved spin
matrices are Hermitian and positive within numerical precision.

\paragraph{Completed central-surface quadrature.}
The refinement fixes source L5 and target R12 at $t=0.60$.
It requires two successive relative changes below $10^{-4}$
for each source integral, target integral, and normalized ratio,
for both signs, both resolved components, and the spin sum.
The source stabilizes at N7; the target continues to N10.
\begin{center}
\begin{tabular}{rrrrr}
\toprule
Source $N$ & Target $N$ & $\mathcal R(s=+1)$ & $\mathcal R(s=-1)$ & $\mathcal R(\text{sum})$\\
\midrule
3 & 3 & 0.252925745 & 0.252836380 & 0.252879225\\
4 & 4 & 0.252306619 & 0.252134476 & 0.252217026\\
5 & 5 & 0.250782696 & 0.250562482 & 0.250668084\\
6 & 6 & 0.250255801 & 0.250023212 & 0.250134747\\
7 & 7 & 0.250108170 & 0.249871433 & 0.249984956\\
7 & 8 & 0.250068188 & 0.249829678 & 0.249944051\\
7 & 9 & 0.250056418 & 0.249817280 & 0.249931954\\
7 & 10 & 0.250052595 & 0.249813244 & 0.249928019\\
\bottomrule
\end{tabular}
\end{center}
The two final maximum source changes are
$9.1214\times10^{-6}$ and $9.1781\times10^{-6}$;
the corresponding target changes are $4.9630\times10^{-5}$
and $1.6157\times10^{-5}$. The full complex matrices also pass
two successive relative Frobenius changes below $10^{-4}$:
source $6.4870\times10^{-6}$, $8.7186\times10^{-6}$;
target $4.8452\times10^{-5}$, $1.5755\times10^{-5}$.
The interrupted partial source N8 grid is excluded.

The final source and target spin sums are
$4.018798024166\times10^{-10}$ and
$1.315213852318\times10^{-9}$.
The resolved-component ratios are $0.249767844257$ and
$0.250092639457$. The difference between the arguments of the
off-diagonal integrated matrix entries is
$-4.802614953\times10^{-4}$ radians; the Frobenius distance
between the trace-normalized matrices is
$5.706858874\times10^{-4}$.
Thus the local normalization still fails the proposed global
comparison. Quadrature convergence alone does not establish the
normalization, interacting spin transport, or block remainder.

\subsection{The conditional factor-four experiment}

The next saved study fixes the explicit assumption
\begin{equation}
 M_{\mathrm{provisional}}=4M_{\mathrm{local}}.
 \label{eq:partition-factor-four}
\end{equation}
This replaces $B_{L,\eta}B_{R,\eta'}/8$ by
$B_{L,\eta}B_{R,\eta'}/2$ on the same supported sectors and
multiplies every matrix entry, including interference, once.
The constants, chiral blocks, primary factors, spin phases,
and spectral measure remain the same. Consequently the
fixed-L5 quadrature ratios above are multiplied by exactly
four; the central final values become
$1.0002103786$, $0.9992529741$, and $0.9997120778$ for the
two signs and their sum. The local identity/torus test does
not derive this additional genus-two normalization.

\paragraph{Source order and independent quadrature.}
For a separate block-order study, complete complex coefficient
banks through total level eight are retained on two source
grids: the archived adapted N7 rule (343 nodes) and a panel
N12 rule (1,728 nodes). Each node retains eight channel/parity
tables with 285 exponent triples through total level eight.
All 2,071 nodes are reused at the five saved plumbing points.
The panel rule has 2, 8, and 2 nodes on the intervals
$[0,0.1]$, $[0.1,2.5]$, and $[2.5,\infty)$, respectively.
For exact reproduction, its frozen momentum nodes and weights
are the \texttt{nodes} entries in
\path{Data Set/nsrr_per_edge8_generic04_20260914/bundle/config.json};
the adapted rule is stored in
\path{Data Set/nsrr_moduli_extension_20260913/coefficient_bank/N7/manifest.json}.
The bank originally built for another surface is independent
of $q$; the comparison reevaluates its coefficients at
\eqref{eq:partition-surfaces} and does not import that older
surface's scalar contraction or primary factors.

At the central surface, with panel N12 for the source and
fixed N10/R12 for the target, the conditional results are
\begin{center}
\begin{tabular}{rrrrr}
\toprule
Source $L$ & $\mathcal R(s=+1)$ & $\mathcal R(s=-1)$ & $\mathcal R(\text{sum})$
 & Nodal change\\
\midrule
3 & 0.9999375694 & 0.9991169498 & 0.9995104611 & ---\\
4 & 1.0002406278 & 0.9992798038 & 0.9997405473 & $4.177\times10^{-4}$\\
5 & 1.0002014989 & 0.9992453421 & 0.9997038475 & $5.896\times10^{-5}$\\
6 & 1.0002080669 & 0.9992517683 & 0.9997103417 & $8.945\times10^{-6}$\\
7 & 1.0002068305 & 0.9992506908 & 0.9997091880 & $1.503\times10^{-6}$\\
8 & 1.0002070665 & 0.9992508688 & 0.9997093938 & $2.668\times10^{-7}$\\
\bottomrule
\end{tabular}
\end{center}
The last column is the largest, over the five scalar
observables, of the sum of absolute changes of the weighted
node contributions divided by the higher-order integral.
Absolute values are taken before summing nodes, so cancellation
between momenta cannot hide the measured order change. It is
a finite-difference diagnostic, not a rigorous tail bound.
Across all five surfaces, the largest such L7-to-L8 change
is $9.247\times10^{-7}$.

At L5, the largest scalar relative difference of panel N12
from the completed Laguerre N7 source is $8.878\times10^{-6}$;
from adapted N7 it is $5.705\times10^{-6}$.
These are independent source-grid controls; they do not replace
the target quadrature refinement.

\paragraph{Target order and validation of saved banks.}
The target controls, both on the N3 grid, are
\begin{center}
\begin{tabular}{rrr}
\toprule
Target recursion step & Largest scalar relative change & Matrix relative change\\
\midrule
R8 to R12 & $3.915\times10^{-7}$ & $2.388\times10^{-7}$\\
R12 to R16 & $1.206\times10^{-9}$ & $7.852\times10^{-10}$\\
\bottomrule
\end{tabular}
\end{center}
The fresh R16 control contains all 27 N3 nodes. There is no
R16 recomputation of the N10 target integral in this record.

Three panel nodes (head, bulk, and tail) were independently
recomputed through L8 with the then-current native engine,
using all four equal-$\eta$ channels and evaluating $f=1$
directly rather than reconstructing it by parity. The maximum
scaled coefficient discrepancy is $3.308\times10^{-34}$;
the projected complex blocks agree at saved floating-point
precision. Independent 40-digit evaluation of the supplied
three-point constants agrees to $4.108\times10^{-14}$ relatively.
This validates selected nodes, not a fresh recomputation of
the entire bank.

Nine unit tests check independent complex antichiral inputs,
all tube signs, the single application of the factor four,
diagonal and interference terms, phase preservation, unchanged
descendant blocks and primary propagation at the resummed API
boundary, rejection of the new normalization option for the
legacy matrix, and the local BPZ/primary-bookkeeping regressions.
Their saved commands, hashes, and outcomes are in
\texttt{tests.json} and \texttt{bank\_validation.json}.

\paragraph{Final conditional result.}
At source L8/panel N12 and target R12/Laguerre N10:
\begin{center}
\begin{tabular}{lrr}
\toprule
Observable & $\mathcal R$ & $100(\mathcal R-1)$\\
\midrule
$s=+1$ & 1.0002070665 & $+0.020707\%$\\
$s=-1$ & 0.9992508688 & $-0.074913\%$\\
Resolved 00 & 0.9990667801 & $-0.093322\%$\\
Resolved 11 & 1.0003698399 & $+0.036984\%$\\
Spin sum & 0.9997093938 & $-0.029061\%$\\
\bottomrule
\end{tabular}
\end{center}
The conditional source spin sum is
$1.607514893908\times10^{-9}$; the target is
$1.315213852318\times10^{-9}$ before the free-frame correction.
The off-diagonal phase difference is
$-4.796613033\times10^{-4}$ radians, and the
trace-normalized matrix distance is $5.715417354\times10^{-4}$.
A positive scalar multiplier does not change phases; these
matrix phases are also distinct from an overall chiral
Pfaffian phase. The reported physical-slice scalar integrals
are positive.

The factor four removes the gross discrepancy but leaves
nonzero residuals. Only the central surface has refined
target quadrature; the other four retain coarse N3 targets.
Finite changes in block order and quadrature are not certified
error bounds. Neither this conditional agreement nor the
independent local checks prove the extra normalization or the
interacting spin transport.

\subsection{Reproduction and the authoritative saved records}

The local derivation and numerical tests are implemented in
\path{Code/genus_2/nsrr_bilinear_sewing.py},
\path{Code/genus_2/nsrr_bilinear_state_oracle.py}, and
\path{Code/genus_2/audit_nsrr_bilinear_sewing.py}.
The oracle contracts full physical Ward/BPW states without
using the reduced block matrix. The audit driver writes Clifford identities, two-point comparisons,
descendant checks, torus limits, free-theory phases, and input
hashes to \path{Data Set/nsrr_bilinear_sewing_20260915/}.
That original output directory is absent from the current checkout;
its local-test counts and errors above are those reported in the
historical note, not a fresh audit of missing raw records.
The three cross-channel/convergence directories listed at the
start of this appendix are present and retain their results.

The following are the archived driver commands, run from the
repository root with Python 3, NumPy, SciPy, mpmath, and the
native Ramond executable available. Use a matching code/data
snapshot: the drivers validate stored configuration and
implementation hashes, so an old result must not be silently
relabeled as a run of a newer executable.
\begin{verbatim}
export OPENBLAS_NUM_THREADS=1 OMP_NUM_THREADS=1
export PYTHONDONTWRITEBYTECODE=1
python3 Code/genus_2/audit_nsrr_bilinear_sewing.py

python3 Code/genus_2/test_nsrr_bilinear_cross_channel.py --workers 3
python3 Code/genus_2/check_cross_channel_source_blocks.py
python3 Code/genus_2/report_nsrr_bilinear_cross_channel.py

python3 Code/genus_2/reproduce_nsrr_bilinear_quadrature.py --workers 6
python3 Code/genus_2/audit_nsrr_bilinear_quadrature.py
python3 Code/genus_2/check_nsrr_quadrature_matrices.py
python3 Code/genus_2/report_nsrr_quadrature_limit.py

python3 Code/genus_2/check_nsrr_factor4_convergence.py \
  --stage all --workers 4
\end{verbatim}
These drivers reuse completed matching records and may compute
missing nodes; they are not all inexpensive report generators.
The conditional driver consumes the complete coefficient banks
listed above and saves its fresh target control under
\texttt{target\_R16\_N3/}. A pointwise resummed worker selects
the new local matrix and the conditional multiplier with
\begin{verbatim}
--sewing-convention human-bilinear
--normalization provisional-times-four
--physical-lifts-slots 1 -1 1
\end{verbatim}
The other tested sign uses \texttt{-1 -1 1}. The older option
\texttt{legacy-times-four} refers to a different interference
matrix and is not this experiment. Likewise, the frozen
pointwise worker's older target scalar is not the full
transported-spin comparison above.

The cross-channel directory's \texttt{config.json} fixes
geometry, spin transport, logarithms, cutoffs, and hashes;
\texttt{source.json}, \texttt{target\_summary.json}, and the
node records retain the complex data. Its
\texttt{comparison.json}, \texttt{result.json}, and
\texttt{fresh\_source\_checks.json} record the initial ratios
and spot checks. The completed quadrature result is
\texttt{final\_result.json}, supplemented by
\texttt{target\_refinement.csv}, \texttt{matrix\_convergence.json},
and \texttt{final\_provenance.json}. Its older
\texttt{summary.json} stops at matched N7 grids and is not
the final N7/N10 conclusion.

In the conditional directory, \path{result.json},
\path{quadrature.csv}, \path{block_orders.csv}, and
\path{source_banks.json} retain the ratios and complex matrices.
Input SHA-256 hashes are in \path{provenance.json} and
\path{source_provenance.json}.
\texttt{bank\_validation.json}, \texttt{verification.json},
and \texttt{tests.json} distinguish saved-node validation
from complete integration. Earlier scalar-only trial
contractions and their marked-spin ratios cannot reconstruct
the new matrix and are not additional crossing evidence.

A lightweight audit of the final unscaled ratios requires no
block evaluations or momentum integrations:
\begin{verbatim}
python3 C++/tools/section6_partition_audit.py \
  --output /tmp/section6-partition-audit.json
\end{verbatim}
It reads the completed N7/N10 record, recomputes the ratios
from the two integrals and free-frame factor, and records
the assumed nature of the multiplier explicitly. It does
not replace the conditional L8 study or establish crossing.

\subsection{Resolution from the Liouville identity residue}
\label{app:partition-normalization-resolution}

The September 24, 2026 audit resolves the factor four without fitting
the genus-two integral. The raw Upsilon coefficients in
\eqref{eq:partition-constants} multiply the normalized ordered forms
as $E,O$, not $E/2,O/2$. Their physical Ramond-family ground values
are consequently $E+O,E-O$. The two-point normalization fixes this
choice independently of any plumbing geometry.

To see this, take the external NS Liouville charge to be
$\alpha_3=\varepsilon\to0^+$, so
$p_3=i(Q/2-\varepsilon)$, and set $p_1=p_2=p>0$.
Divide the NS--R--R coefficients by the all-NS coefficient $C$ at
these same momenta. The common external NS leg and cosmological
factor cancel; only the two internal-leg normalizations differ.
The positive-momentum gamma metrics in \eqref{eq:partition-legs}
are exactly $1/b$ in NS and $1$ in R. Therefore
\begin{align}
 \frac{E}{C}&=
 \frac{\mathcal N_{\R}(p)^2}{b^2\mathcal N_{\NS}(p)^2}
 \frac{\Upsilon_{\NS}(\varepsilon+2ip)
       \Upsilon_{\NS}(Q-\varepsilon+2ip)}
      {\Upsilon_{\R}(\varepsilon+2ip)
       \Upsilon_{\R}(Q-\varepsilon+2ip)}
 \ \longrightarrow\ 1,\nonumber\\
 \frac{O}{C}&=
 \frac{\mathcal N_{\R}(p)^2}{b^2\mathcal N_{\NS}(p)^2}
 \frac{\Upsilon_{\NS}(\varepsilon)^2}
      {\Upsilon_{\R}(\varepsilon)^2}
 \ \longrightarrow\ 0.
 \label{eq:partition-identity-residue}
\end{align}
The first limit uses $\Upsilon_s(Q-x)=\Upsilon_s(x)$ and
\eqref{eq:partition-legs}; the second uses the simple zero of
$\Upsilon_{\NS}$ at zero. The same limits hold when $p_2-p_1$ is
of order $\varepsilon$. The singular factors
$\Upsilon_{\NS}(\varepsilon\pm i(p_2-p_1))$ are shared by $E$
and $C$, so this compares the residues of the same delta-function
peak. Both $E+O$ and $E-O$ reproduce the NS identity residue, as
required by the equal NS/R metrics. Their former half-sums do not.

This agrees with the equal-metric conventions in
\href{https://arxiv.org/pdf/hep-th/9607120}{Poghossian,
Eqs.~(14)--(15), (51), and (56)}. The historical adapter applied
the half-sum dictionary of
\href{https://arxiv.org/html/2201.05621\#S3.SS1}{Balthazar--Rodriguez--Yin,
Eq.~(3.8)} directly to our raw Upsilon terms; that combination fails
\eqref{eq:partition-identity-residue}. The present convention is
specified by the identity residue, not by a change of Ramond phases.
The earlier module test prescribed $E=2,O=0$ by hand and hence could
not detect this incompatibility with the actual structure constants.

The finite Clifford reduction remains unchanged. Its matrix is
$\tfrac12\bigl(\begin{smallmatrix}1&s\\s&1\end{smallmatrix}\bigr)$
times the product of the ordered-form coefficients. Thus the
denominator eight in the old matrix becomes two, as in
\eqref{eq:app-nonchiral-matrix}. No chiral coefficient, Ward identity,
Gram matrix, or momentum measure changes. A denominator eight is
valid only if its inputs are the family sums $2E,2O$.

The actual special-function limits were checked at $b=0.8,1,1.4$,
$p=0.23,0.67$, and $\varepsilon=10^{-2},10^{-4},10^{-6}$.
At the smallest regulator, offsets
$(p_2-p_1)/\varepsilon=-1,0,1$ were included. All corrected
family ratios there differ from one by less than $10^{-5}$.
At $b=1$ an independent Barnes-$G$ implementation agrees within
$5.47\times10^{-17}$, limited by its binary64 return values.
The native Clifford checks pass all 64 cases in machine and
40-digit arithmetic, with maximal scaled errors
$3.15\times10^{-15}$ and $3.25\times10^{-40}$, respectively.
They include the identity trace two with $E=1,O=0$ after dividing
out the common NS identity residue.

A fresh native calculation at the central saved surface gives
\begin{center}
\begin{tabular}{lc}
\toprule
observable & corrected ratio $\mathcal R$\\
\midrule
$s=+1$ & $1.000006186773435$\\
$s=-1$ & $1.000005527870852$\\
resolved $00$ & $1.000007836181439$\\
resolved $11$ & $1.000003799186142$\\
spin sum & $1.000005843991480$\\
\bottomrule
\end{tabular}
\end{center}
The source uses $N=7$ and total descendant level 5, the target
$N=10$ and total level 8, both with 40-digit arithmetic and four
workers. They took $34.56$ and $54.02$ seconds, excluding compilation
and input preparation. The maximal residual is $7.84\times10^{-6}$;
the source level-4-to-5 shift is at most $5.90\times10^{-5}$ and
the target level-7-to-8 shift at most $8.62\times10^{-8}$.
These shifts are convergence diagnostics, not rigorous error bounds.
The momentum rule was retained without a new refinement study.

The native target also corrects a separate Schottky spin-lift
assignment in the historical pointwise evaluator: the half-multiplier
signs are fixed coefficientwise before complex plumbing evaluation.
Its independent low-level PBW check is recorded in
\path{C++/results/partition_conventions_2026-09-24/}.
The new identity tests, all fresh momentum nodes, and integrated
results are in
\path{C++/results/partition_normalization_resolution_2026-09-24/}.
The derivation and reproduction commands are in
\path{C++/PARTITION_NORMALIZATION.md} and \path{C++/PARTITION.md}.
The native reducer rejects output with the old coefficient
normalization. The comparison supports the specified interacting
spin dictionary at this surface; it does not prove it for arbitrary
geometry.

\endgroup
